\section{The local logarithmic square}
\label{sec:local-square}

The numerical invariant used below can be computed on a partial resolution.
We first specify its local intersection convention, then express it as a
difference of ordinary projective intersection numbers.

Let $(V,p)$ be a normal complex surface singularity and let
$\rho:\widetilde V\to V$ be a good resolution: $\widetilde V$ is smooth,
$\rho$ is proper and isomorphic away from $p$, and the full reduced
exceptional divisor $E=\sum_jE_j$ has simple normal crossings.
Its intersection matrix is negative definite.  Consequently, for a rational
divisor $Q$ on $\widetilde V$, there is a unique rational exceptional cycle
$\nu(Q)$ satisfying
\[
  \nu(Q)\cdot E_j=Q\cdot E_j\qquad\text{for every }j.
\]
Write the relative Zariski decomposition as
\[
 K_{\widetilde V}+E=P+N,
\]
where $N$ is an effective rational exceptional cycle, $P\cdot E_j\geq0$
for every $j$, and $P\cdot E_j=0$ whenever $E_j$ occurs in $N$.
Wahl's local logarithmic invariant \cite{Wahl1990} is
\begin{equation}\label{eq:beta-definition}
 \beta(V,p)=-\nu(P)^2.
\end{equation}
We also write $P_{\mathrm{loc}}=\nu(P)$ for this numerical exceptional cycle.
These are the relative decomposition and intersection conventions used by
Okuma \cite[Section~2, pp.~649--650]{Okuma2004}.

The value is unchanged on a higher good resolution and agrees with the
value on the minimal semigood resolution, where the exceptional divisor
has normal crossings, allowing irreducible components to have self-nodes.
Indeed, let $\pi:\widehat V\to\widetilde V$ be a further resolution,
isomorphic off $E$, with full reduced exceptional divisor $\widehat E$ over
$p$.  Log canonicity of the normal crossings pair gives
\[
 K_{\widehat V}+\widehat E
   =\pi^*(K_{\widetilde V}+E)+G,
 \qquad G\geq0,
\]
where $G$ is $\pi$-exceptional.  Hence its relative decomposition is
$\pi^*P+(\pi^*N+G)$: the first summand is relatively nef and is orthogonal
to every component of the second.  Uniqueness of the relative
decomposition applies.  Also $\nu(\pi^*P)=\pi^*\nu(P)$, as one checks
against both strict transforms and $\pi$-exceptional curves.  The square
therefore stays unchanged.  A common higher resolution compares any two
choices, including a semigood one.

\begin{lemma}[Projective computation of the local square]
\label{lem:projective-local-square}
Let $V$ be an integral normal projective complex surface, smooth outside
a point $p$, with Cartier canonical divisor.  Let $b:S\to V$ be a
projective birational morphism from a normal surface, isomorphic outside
$p$.  Suppose that the reduced inverse image of $p$ is a divisor $B$,
that $(S,B)$ is log canonical, and that $K_S+B$ is $b$-nef.  Then
\begin{equation}\label{eq:projective-local-square}
 \beta(V,p)=K_V^2-(K_S+B)^2.
\end{equation}
\end{lemma}

\begin{proof}
Choose a log resolution $\pi:Z\to S$, isomorphic outside $B$, and put
$\rho=b\pi$.  Let $E$ be the full reduced exceptional divisor of $\rho$.
With compatible canonical divisors, log canonicity gives
\[
 K_Z+E=\pi^*(K_S+B)+N,
\]
where $N$ is effective and $\pi$-exceptional.  The pullback on the right
is $\rho$-nef and has intersection zero with every component of $N$.
It is therefore the positive part $P$ of the relative decomposition.

The divisor $Z_P=P-\rho^*K_V$ is exceptional: the two summands agree
outside $\rho^{-1}(p)$.  Moreover, $\rho^*K_V$ has intersection zero with
every exceptional curve, so $Z_P=\nu(P)$.  The projection formula gives
\[
 P^2=K_V^2+Z_P^2=(K_S+B)^2.
\]
Together with \eqref{eq:beta-definition}, this proves the formula.
\end{proof}

\begin{lemma}[Persistence of log canonicity]\label{lem:lc-persistence}
In a holomorphic family of isolated hypersurface singularities in three
spatial variables with constant Milnor number over a connected smooth
complex curve, log
canonicity of the normal surface germ is constant.
\end{lemma}
\begin{proof}
By Varchenko's theorem \cite{Varchenko1982}, as recovered by
Steenbrink, the spectrum of an isolated hypersurface singularity
is constant in a $\mu$-constant deformation
\cite[Theorem~2.8]{Steenbrink1985}.  Its smallest number plus one is the
complex singularity index \cite[Theorem~2.11]{Steenbrink1985}, and the
log canonical threshold is the minimum of this index and one
\cite[Theorem~9.5]{Kollar1997Pairs}.  For a normal hypersurface in a smooth
ambient space, adjunction and inversion of adjunction identify log
canonicity with threshold one \cite{Kawakita2007}.  To apply the algebraic
inversion theorem to an individual analytic germ, replace its defining
function by a sufficiently high Taylor polynomial using
Lemma~\ref{lem:individual-jets}; analytic right equivalence preserves both
the threshold and log canonicity.  Thus log canonicity, and its failure,
is constant in the family.  This use of the spectrum concerns the
individual isolated function germs; it does not assume topological
triviality of the family.
\end{proof}
